\subsection{The Common Log} 
\begin{definition}[Common Logarithm]
When we write the symbol $\log$ without a base (that is, ``$\log$'' rather than ``$\log_b$''), it refers to the \term{Common Logarithm}, which in our context is the logarithm base 10.
\end{definition}
\begin{Example}
Rewrite $\log 100=2$ in exponential form, and determine whether it is true.
\begin{lpic}{}{1}
\end{lpic}
Since the exponential statement is \makeblank[1in], the log form statement is also \makeblank[1in].
\end{Example}
One nice thing about the common log is that the calculator has a button for it.
\begin{Example}
Use your calculator to compute each logarithm. Round to three decimal places if necessary.
\begin{itemize}
\item $\log 1000=\makeblanksmall$ (Check: \makeblank)
\item $\log 0.01=\makeblanksmall$ (Check: \makeblank)
\item $\log 17=\makeblanksmall$ (Check: \makeblank)
\item $\log 1945=\makeblanksmall$(Check: \makeblank)
\end{itemize}
\end{Example}
Since our number system is \makeblank, the common log can be approximated by counting the \makeblank[1in] of the number. If a number has $n$ digits to the left of the decimal point, its common log will be between \makeblanksmall and \mbox{\makeblanksmall.} If its first nonzero digit is $n$ places to the right of the decimal point, its common log will be between \makeblanksmall and \makeblanksmall.
\begin{Example}
Estimate each common log by counting the digits. Then find the log on your calculator, rounded to three decimal places.
\begin{itemize}
\item $\log 19271$
\begin{itemize}
\item Counting: $\makeblanksmall<\log 19271<\makeblanksmall$
\item Calculator: $\log 19271\approx\makeblanksmall$
\end{itemize}
\item $\log 0.0083$
\begin{itemize}
\item Counting: $\makeblanksmall<\log 0.0083<\makeblanksmall$
\item Calculator: $\log 0.0083\approx\makeblanksmall$
\end{itemize}
%\item $\log 1.24$
%\begin{itemize}
%\item Counting: $\makeblanksmall<\log 425<\makeblanksmall$
%\item Calculator: $\log 425\approx\makeblanksmall$
%\end{itemize}
\end{itemize}
\end{Example}